% ============================================================
%  ALDIMI Predict -- Hito 4 (Trabajo Final)
%  Secciones nuevas para completar el informe final segun el
%  enunciado (seccion 9: Estructura del Informe Final).
%  Integrar a TB1_ALDIMI_v5.tex (o crear TF_ALDIMI_v6.tex) y
%  recompilar con XeLaTeX (2 pasadas).
%  Los marcadores [COMPLETAR: ...] requieren datos del equipo.
% ============================================================

% ──────────────────────────────────────────────────────────
% I. RESUMEN EJECUTIVO  (colocar al inicio, antes de la sec. 1)
% ──────────────────────────────────────────────────────────
\section*{Resumen Ejecutivo}
\addcontentsline{toc}{section}{Resumen Ejecutivo}

El Albergue Divina Misericordia (ALDIMI) atiende gratuitamente a niños y
adolescentes con cáncer en situación de extrema pobreza y se encuentra en
plena expansión (\emph{ALDIMI 2.0}): duplicará su capacidad de 50 a 100
familias. Esta expansión es insostenible con procesos manuales: el kardex de
almacén en hojas de cálculo no permite anticipar quiebres de stock, y la
priorización de la atención de pacientes depende exclusivamente de la
observación del personal.

\textbf{ALDIMI Predict} aborda ambos problemas con dos frentes de Machine
Learning desarrollados bajo la metodología CRISP-DM:

\begin{itemize}[leftmargin=1.6em]
  \item \textbf{Frente 1 --- Predicción de inventario:} modelos de
        clasificación que anticipan el quiebre de stock de cada artículo a
        7 y 14 días. El modelo seleccionado a 7 días (XGBoost) alcanza
        \textbf{F1 = 0.9865} y \textbf{AUC = 0.9982}; a 14 días
        (Random Forest), F1 = 0.9721 y AUC = 0.9901.
  \item \textbf{Frente 2 --- Clasificación de riesgo de pacientes:} un
        clasificador multiclase (XGBoost) que estima el nivel de prioridad
        de atención (Bajo/Medio/Alto) con \textbf{F1-Macro = 0.75} y
        \textbf{AUC-Macro = 0.9164} sobre 600 historias sintéticas derivadas
        de las fichas reales de ingreso y reingreso del albergue.
\end{itemize}

Ambos modelos se despliegan en un \emph{dashboard} interactivo (Streamlit)
diseñado para los directivos del albergue: alertas de stock accionables,
evaluación individual de pacientes con factores de riesgo explicados
(valores SHAP) y indicadores globales de gestión. El valor predictivo
aportado se traduce en compras y campañas de donación planificadas con una
o dos semanas de anticipación, y en la priorización temprana de pacientes
con alta probabilidad de complicaciones, en línea con los ODS~2, 3 y 10.

% ──────────────────────────────────────────────────────────
% FASE 6 CRISP-DM: DESPLIEGUE  (nueva sección, después de la
% evaluación / antes de la sección de valor del proyecto)
% ──────────────────────────────────────────────────────────
\section{Fase 6 --- Despliegue (Deployment) e Implementación MLOps}

\subsection{Arquitectura del pipeline}

El sistema sigue un pipeline de cuatro etapas, del dato crudo a la decisión:

\begin{enumerate}[leftmargin=1.8em]
  \item \textbf{Ingesta:} el kardex \texttt{ALMAcen upc.xlsx} y las fichas
        de ingreso/reingreso de pacientes se transforman en los datasets
        procesados (\texttt{data/processed/}): consolidado por artículo,
        serie semanal de inventario y dataset sintético de pacientes.
  \item \textbf{Entrenamiento:} el notebook
        \texttt{notebooks/aldimi\_analisis\_modelado.ipynb} ejecuta la
        preparación, el entrenamiento con \texttt{RandomizedSearchCV}
        (50--60 iteraciones) y validación cruzada estratificada (5 folds),
        y serializa los artefactos con \texttt{joblib}.
  \item \textbf{Artefactos versionados:} \texttt{models/models\_inventario.pkl}
        (XGBoost 7d + Random Forest 14d, con \emph{label encoder} y
        \emph{scaler}) y \texttt{models/models\_pacientes.pkl} (XGBoost
        multiclase con las 45 features de entrenamiento). Empaquetar el
        preprocesamiento junto al modelo garantiza que la inferencia use
        exactamente las mismas transformaciones que el entrenamiento.
  \item \textbf{Serving:} \texttt{dashboard.py} (Streamlit) carga los
        artefactos con caché (\texttt{st.cache\_resource}) y expone la
        inferencia en tiempo real a través de formularios; se ejecuta en
        local y se publica para demostraciones mediante t\'unel ngrok
        (\texttt{lanzar\_demo.py}).
\end{enumerate}

\subsection{Estrategia de mantenimiento y reentrenamiento}

\begin{itemize}[leftmargin=1.6em]
  \item \textbf{Reentrenamiento periódico:} al cierre de cada mes (o al
        incorporarse datos nuevos del módulo de IA), se re-ejecuta el
        notebook y se regeneran los \texttt{.pkl}; el dashboard los toma
        automáticamente sin cambios de código.
  \item \textbf{Monitoreo de deriva:} comparar la distribución semanal de
        \texttt{salidas\_semana} y la ocupación real contra los rangos de
        entrenamiento (50--100 familias); si la ocupación supera el rango,
        el modelo debe reentrenarse antes de confiar en sus alertas.
  \item \textbf{Criterio de aceptación:} mantener F1 $\geq$ 0.95 en
        inventario y F1-Macro $\geq$ 0.70 en pacientes sobre el conjunto de
        prueba; por debajo de estos umbrales las alertas pasan a modo
        consultivo (revisión manual obligatoria).
\end{itemize}

\subsection{Manual de usuario}

Se elaboró un manual de usuario orientado al personal no técnico del
albergue (instalación, navegación, interpretación de alertas y de niveles de
riesgo, y solución de problemas). Se adjunta en
\texttt{docs/manual\_usuario\_ALDIMI.docx} dentro del paquete de entrega.

% ──────────────────────────────────────────────────────────
% IV (complemento). ÉTICA DE DATOS Y ANÁLISIS DE SESGOS
% ──────────────────────────────────────────────────────────
\section{Ética de datos y análisis de sesgos}

\subsection{Protección de datos y anonimización}

Los datos de pacientes utilizados para el entrenamiento son
\textbf{sintéticos}: se generaron respetando las distribuciones y reglas
clínicas de las fichas reales de ingreso y reingreso de ALDIMI, pero sin
contener ningún dato identificable de un niño real. Los identificadores
(\texttt{id\_paciente}) son códigos anónimos. Esto permite desarrollar y
compartir el sistema sin exponer información de salud de menores, en
cumplimiento de la Ley de Protección de Datos Personales (Ley N.º 29733).

\subsection{Sesgos identificados y mitigaciones}

\begin{itemize}[leftmargin=1.6em]
  \item \textbf{Desbalanceo de clases:} la clase ``Alto'' es minoritaria.
        Se mitigó con \texttt{class\_weight} balanceado (Random Forest) y
        \texttt{scale\_pos\_weight} (XGBoost), y se evaluó con F1-Macro
        --- que pondera por igual las tres clases --- en lugar de accuracy.
  \item \textbf{Sesgo socioeconómico:} variables como el grado de
        instrucción del cuidador o la procedencia podrían inducir a que el
        modelo asigne sistemáticamente mayor riesgo a familias de menor
        nivel educativo o de zonas rurales. El análisis SHAP muestra que
        las variables dominantes son clínicas (etapa del cáncer, estado
        físico y nutricional, infección); las variables sociales actúan
        como moduladores de contexto. Aun así, se recomienda auditar
        periódicamente las tasas de clasificación ``Alto'' por procedencia
        y nivel educativo para detectar disparidades.
  \item \textbf{Automatización de decisiones clínicas:} el sistema se
        diseñó como \emph{apoyo} a la priorización, nunca como decisor
        final. Toda clasificación ``Alto'' exige confirmación del personal
        de salud (human-in-the-loop), y el dashboard muestra los factores
        detectados para que la recomendación sea auditable y explicable.
  \item \textbf{Datos sintéticos:} el uso de datos sintéticos evita riesgos
        de privacidad pero puede no capturar toda la variabilidad real. En
        la fase de integración con datos reales (módulo de IA) los modelos
        deberán revalidarse antes de su uso operativo.
\end{itemize}

% ──────────────────────────────────────────────────────────
% V. CONFLUENCIA CON EL ECOSISTEMA "ALDIMI CORE AI"
% ──────────────────────────────────────────────────────────
\section{Confluencia con el ecosistema ALDIMI Core AI}

\subsection{Contrato de datos con el módulo de IA}

El curso de IA (1ASI0404) desarrolla el módulo de captura: OCR de fichas y
recetas, y chatbot de registro. El punto de encuentro es la \textbf{base de
datos común}: los campos que el OCR extrae de las fichas de ingreso
(edad, diagnóstico, etapa, estado nutricional, etc.) coinciden con las
variables de entrada del clasificador de riesgo, y los registros de
movimientos de almacén alimentan la serie semanal de inventario.

[COMPLETAR: describir la implementación real acordada con el grupo de IA ---
motor de BD (p.~ej. PostgreSQL/Supabase/SQLite), esquema de tablas, y
mecanismo de lectura (consulta directa / API REST / exportación CSV
programada). Adjuntar captura o extracto del esquema como evidencia.]

\subsection{Flujo end-to-end}

El flujo demostrado en el video de impacto es: (1) el OCR del módulo de IA
captura una ficha de ingreso y la persiste en la BD común; (2) el pipeline
de ML lee el nuevo registro, aplica el preprocesamiento empaquetado en los
artefactos y genera la predicción; (3) el dashboard muestra la alerta de
stock o el nivel de riesgo del paciente para la toma de decisiones.

\subsection{Lecciones aprendidas de la integración}

\begin{itemize}[leftmargin=1.6em]
  \item Definir el \emph{contrato de interfaz} (nombres, tipos y valores
        v\'alidos de los campos) al inicio evita retrabajos: los valores
        categ\'oricos del OCR deben coincidir exactamente con las categor\'ias
        vistas en el entrenamiento.
  \item Empaquetar \emph{scaler} y \emph{encoders} dentro del \texttt{.pkl}
        elimina inconsistencias entre los equipos.
  \item [COMPLETAR: 1--2 lecciones específicas del trabajo real con el
        grupo de IA.]
\end{itemize}

% ──────────────────────────────────────────────────────────
% VI. ANEXOS  (colocar al final, antes de las referencias)
% ──────────────────────────────────────────────────────────
\section*{Anexos}
\addcontentsline{toc}{section}{Anexos}

\begin{itemize}[leftmargin=1.6em]
  \item \textbf{Repositorio de código:}
        [COMPLETAR: URL de GitHub/GitLab con el código limpio y comentado].
  \item \textbf{Diccionario de datos:} archivo
        \texttt{docs/diccionario\_datos\_ALDIMI.xlsx} --- describe las
        variables, tipos, rangos y porcentaje de nulos de los tres datasets
        (\texttt{dataset\_semanal}, \texttt{dataset\_completo},
        \texttt{pacientes\_sintetico}).
  \item \textbf{Manual de usuario:}
        \texttt{docs/manual\_usuario\_ALDIMI.docx}.
  \item \textbf{Notebook de entrenamiento:}
        \texttt{codigo/notebooks/aldimi\_analisis\_modelado.ipynb}
        (reproduce el pipeline completo de CRISP-DM fases 3--5).
  \item \textbf{Video de impacto:} [COMPLETAR: enlace al video de
        demostración end-to-end integrado con el curso de IA].
\end{itemize}
